\begin{abstract}
Single-step adversarial training is efficient but can abruptly lose robustness
through catastrophic overfitting (CO). We study an online safety trigger based
on the cosine consistency of two input gradients obtained within an FGSM
parameter update. The CIFAR-10 threshold is selected on seeds 17, 23, and 42
and frozen before held-out evaluation. It detects 7/7 held-out
CIFAR-10/PreActResNet-18 events within 60 updates. With the threshold unchanged,
it detects 4/5 CIFAR-10/ResNet-18 events, including 2/5 within 60 updates. After
one dataset-specific calibration, it detects both observed held-out CIFAR-100
events; the three CIFAR-100 non-event trajectories contain one alarm episode.
Eight randomized-FGSM runs at $\epsilon=8/255$ remain stable and produce no
alarms. At $\epsilon=16/255$, all three stress-test runs naturally collapse;
the frozen CIFAR-10 detector warns 20--60 updates beforehand, and the tested
\emph{immediate-replay} ablation raises terminal PGD-10 accuracy from 0 to
$26.48\pm0.83\%$. In a separate five-seed timing audit, both
\emph{immediate-switch} and the rollback ablation \emph{immediate-replay}
start PGD-2 before CO in 5/5 cases, versus 3/5 for
next-epoch escalation. For immediate-switch, final PGD-10 accuracy is
$47.05\pm0.74\%$; its $+0.22$-point difference from immediate-replay has a paired 95\%
interval crossing zero. Immediate switching provides the timing benefit; rollback
adds no reliable endpoint gain. In controlled reimplementations, GradAlign
achieves higher robustness, while the proposed policy uses lower measured core
training time. For the fixed seed-101 immediate-switch checkpoint, full-test PGD-50,
five-restart PGD-20, and fixed-subset AutoAttack-standard accuracy are 45.52\%,
45.66\%, and 42.30\%, respectively. Gradient consistency is used as an
empirical warning signal; conditional PGD-2 supplies the recovery.
\end{abstract}
